\section{Công nghệ và quyết định lựa chọn}

Chương này trình bày các nhóm công nghệ được lựa chọn để xây dựng trợ lý mua sắm
đa phương thức. Nội dung tập trung vào vai trò trong hệ thống, lý do lựa chọn và
các giới hạn cần kiểm soát; các nguyên lý embedding, truy hồi, RAG và các cấu hình
chi tiết được trình bày ở Chương 3, Chương 5 và phụ lục.

\subsection{Tiêu chí và tổng quan lựa chọn công nghệ}

Các công nghệ được đánh giá theo năm tiêu chí chính:
\begin{enumerate}[label=\arabic*., leftmargin=1.15cm]
    \item đáp ứng nghiệp vụ thương mại điện tử và tìm kiếm đa phương thức;
    \item phù hợp với mã nguồn và năng lực triển khai của đề tài;
    \item có hợp đồng giao tiếp rõ ràng, thuận tiện kiểm thử và thay thế;
    \item vận hành được trong môi trường container;
    \item cân bằng chất lượng, độ trễ, chi phí và khả năng quản trị.
\end{enumerate}

Việc lựa chọn không chỉ dựa trên mức độ phổ biến của công cụ mà còn dựa trên khả
năng tích hợp với các thành phần đã có và khả năng thay thế khi cần.

Kiến trúc công nghệ được tổ chức thành các nhóm. React đảm nhiệm giao diện; Spring
Boot và MySQL xử lý nghiệp vụ cùng dữ liệu chuẩn; FastAPI điều phối AI Service;
LangGraph quản lý workflow hội thoại; Bedrock cung cấp LLM và embedding; Qdrant
phục vụ truy hồi vector; Docker, Google Compute Engine và Cloud Run cung cấp môi
trường triển khai. Các nhóm này giao tiếp qua API và không để trình duyệt truy cập
trực tiếp vào cơ sở dữ liệu, vector database hoặc khóa dịch vụ AI.

\begin{table}[H]
    \centering
    \caption{Tổng quan các nhóm công nghệ được lựa chọn}
    \label{tab:technology_stack_overview}
    \renewcommand{\arraystretch}{1.15}
    \begin{tabularx}{\linewidth}{|p{0.25\linewidth}|X|}
        \hline
        \textbf{Nhóm} & \textbf{Công nghệ chính} \\ \hline
        Ứng dụng và dữ liệu & React, Redux Toolkit, Spring Boot, Spring Security, MySQL \\ \hline
        AI và đa phương thức & Python, FastAPI, LangGraph, Amazon Bedrock, YOLOv26, Pillow \\ \hline
        Truy hồi và trạng thái & Cohere Embed v4, Qdrant, hybrid fusion, reranker, Redis \\ \hline
        Triển khai & Docker, Docker Compose, Google Compute Engine, Cloud Run \\ \hline
    \end{tabularx}
\end{table}

\subsection{Công nghệ ứng dụng và dữ liệu nghiệp vụ}

Frontend sử dụng React 19 để xây dựng giao diện theo component \parencite{reactdocs2024}.
React phù hợp với các màn hình có nhiều trạng thái tương tác như bộ lọc, danh sách
kết quả, giỏ hàng, vùng tải ảnh và hội thoại nhiều lượt. Redux Toolkit quản lý các
trạng thái dùng chung như phiên đăng nhập và giỏ hàng; React Router tổ chức điều
hướng phía client; Axios tập trung việc gọi API, xử lý token, timeout và lỗi. Nhóm
công nghệ này giúp tái sử dụng giao diện giữa tìm kiếm văn bản, hình ảnh và kết hợp.
Mọi dữ liệu nghiệp vụ quan trọng vẫn phải được Backend xác nhận, do Frontend không
được tự quyết định giá, tồn kho hoặc trạng thái đơn hàng.

Backend sử dụng Java 21 và Spring Boot 4.0.6 để xây dựng REST API. Spring Boot cung
cấp hệ sinh thái thống nhất cho controller, validation, transaction, security và
truy cập dữ liệu quan hệ \parencite{springbootdocs2024}. Spring Security cùng OAuth2
Client xử lý xác thực và phân quyền; Spring Data JPA ánh xạ entity; RestClient thực
hiện giao tiếp nội bộ với AI Service. Spring Boot được chọn vì phù hợp với các nghiệp
vụ có ràng buộc giao dịch như tồn kho, khuyến mãi, đơn hàng và thanh toán COD. Hạn
chế là cấu hình và thời gian khởi động lớn hơn một API tối giản; vì vậy suy luận AI
được tách thành một dịch vụ Python độc lập.

MySQL 8.4 là nguồn dữ liệu chuẩn cho tài khoản, sản phẩm, tồn kho, giỏ hàng, khuyến
mãi và đơn hàng. Mô hình quan hệ hỗ trợ khóa ngoại, transaction và các ràng buộc
nghiệp vụ cần thiết. Qdrant chỉ là chỉ mục phục vụ tìm kiếm, có thể tái tạo từ MySQL;
tính chất transaction của InnoDB được mô tả trong tài liệu MySQL chính thức
\parencite{mysqlinnodb2026}. Kết quả AI chỉ cung cấp mã sản phẩm và điểm tương đồng,
sau đó Backend đọc lại giá,
tồn kho và trạng thái hiện hành trước khi trả về cho người dùng.

\subsection{Công nghệ AI và xử lý dữ liệu đa phương thức}

AI Service sử dụng Python 3.11, FastAPI và Pydantic. FastAPI phù hợp với API bất đồng
bộ, sinh OpenAPI từ type hints và kết nối với hệ sinh thái xử lý AI \parencite{fastapidocs2024}.
Pydantic chuẩn hóa request/response; HTTPX thực hiện lời gọi đến Backend và các dịch
vụ mô hình; Pillow giải mã, resize và crop ảnh. Ranh giới REST giúp thay đổi provider
embedding hoặc LLM mà không làm thay đổi giao diện người dùng. Đổi lại, hợp đồng API,
timeout và cơ chế xử lý lỗi giữa hai dịch vụ phải được kiểm soát rõ ràng.

LangGraph điều phối workflow hội thoại có trạng thái dưới dạng graph
\parencite{langgraph2024docs}. Các node đảm nhiệm phân tích truy vấn, định tuyến,
truy hồi, lấy dữ liệu sản phẩm và sinh phản hồi; các điều kiện chuyển trạng thái được
giới hạn trước để bảo đảm hành vi có thể kiểm thử. Redis được dùng ở môi trường Cloud
Run để lưu checkpoint và preference memory ngoài tiến trình, phù hợp với vai trò lưu
trữ trạng thái và dữ liệu trung gian của Redis \parencite{redisdocs2026}. Cách tổ chức này phù hợp
với việc container AI Service có thể được thay thế, nhưng cần giới hạn dung lượng,
thời gian lưu trữ và quyền truy cập bộ nhớ.

LLM được sử dụng qua Amazon Bedrock để phân tích ý định, trích xuất thuộc tính, viết
lại truy vấn và diễn đạt câu trả lời. Dữ kiện về sản phẩm không lấy trực tiếp từ tri
thức tham số của mô hình mà được cung cấp bởi tầng truy hồi và Backend. Dịch vụ managed
giúp không phải vận hành GPU cho LLM, nhưng tạo phụ thuộc vào mạng, quota, chi phí và
provider \parencite{awsbedrockdocs2026}; vì vậy AI Service cần timeout, retry có giới hạn,
fallback và ghi nhận model cùng độ trễ.

YOLOv26s được dùng như một bước tăng cường trong xử lý ảnh. Mô hình được fine-tune
trên dữ liệu thời trang để phát hiện các nhóm phụ kiện, túi, quần áo và giày; bounding
box đủ tin cậy được crop trước khi tạo image embedding. Pillow thực hiện giải mã và
chuẩn hóa ảnh. YOLO không phải điều kiện bắt buộc của Image Search: nếu không phát
hiện được vùng phù hợp, hệ thống dùng ảnh đầy đủ. Cách fallback này cho phép cấu hình
Cloud Run tắt YOLO để giảm kích thước container và yêu cầu tài nguyên.

\subsection{Công nghệ truy hồi, xếp hạng và quản lý trạng thái}

Cohere Embed v4 qua Amazon Bedrock tạo vector cho văn bản và hình ảnh trong cùng
pipeline đa phương thức. Hệ thống duy trì collection riêng cho từng modality, mỗi
point liên kết với mã sản phẩm và các thuộc tính dùng để lọc như danh mục, màu, giá
hoặc trạng thái hoạt động. Embedding managed giúp thống nhất mô hình biểu diễn mà
không phải phục vụ mô hình vision--language cục bộ; hạn chế là độ trễ mạng, quota và
chi phí theo request \parencite{cohereembedv42026}.

Qdrant được chọn làm vector database vì hỗ trợ cosine similarity, HNSW, payload
filtering, API Python và self-host bằng Docker \parencite{qdrant_docs}. Search Service
truy hồi các ứng viên từ nhánh văn bản và hình ảnh, quy về mã sản phẩm rồi áp dụng
metadata cùng ràng buộc nghiệp vụ. Cách triển khai này cho phép kiểm soát dữ liệu và
chi phí ở quy mô catalog hiện tại. Đổi lại, nhóm phát triển phải quản lý volume,
snapshot, backup và khả năng tái lập chỉ mục.

Các nhánh truy hồi được hợp nhất bằng hybrid fusion sau khi chuẩn hóa điểm. Reranker
managed chỉ được bật qua feature flag trên tập ứng viên nhỏ, vì nó có thể làm tăng độ
trễ và chi phí. Cấu hình báo cáo ưu tiên fusion không rerank do kết quả đánh giá hiện
tại chưa chứng minh rerank cải thiện các chỉ số chính; thành phần này vẫn được giữ để
thử nghiệm trong các cấu hình sau.

Redis đồng thời cung cấp nơi lưu checkpoint và trạng thái hội thoại khi AI Service
chạy trên Cloud Run. Trạng thái cần được phân tách theo người dùng và phiên, có TTL,
giới hạn dung lượng và không được chứa dữ liệu nhạy cảm ngoài phạm vi cần thiết. Việc
đưa trạng thái ra ngoài container giúp các instance xử lý tiếp cùng một phiên, trong
khi MySQL vẫn giữ vai trò nguồn chuẩn cho dữ liệu giao dịch.

\subsection{Công nghệ triển khai và tổng hợp quyết định}

Docker đóng gói Frontend, Spring Boot Backend, AI Service, MySQL và Qdrant thành các
image hoặc container có cấu hình xác định. Docker Compose mô tả network, health check,
phụ thuộc dịch vụ và named volume, giúp môi trường phát triển cùng kiểm thử được khởi
tạo nhất quán. Frontend và Backend sử dụng build nhiều giai đoạn; secret và cấu hình
môi trường không được đóng trực tiếp vào image. Volume của MySQL và Qdrant vẫn cần
được sao lưu trước khi cập nhật \parencite{dockerdocs2026,dockercompose2026}.

Trong môi trường Google Cloud, GCE VM chạy Frontend, Spring Boot Backend, MySQL và
Qdrant; AI Service được triển khai riêng trên Cloud Run. Cloud Run kết nối đến các
dịch vụ cần thiết qua mạng riêng; Secret Manager cung cấp token Bedrock, internal API
key và Redis URL. GCE phù hợp với thành phần cần volume và cấu hình mạng ổn định,
trong khi Cloud Run phù hợp với AI Service dạng stateless khi checkpoint được đưa vào
Redis. Đổi lại, mô hình này cần quản lý IAM, logging, quota, private networking và
chi phí trên hai kiểu môi trường \parencite{gce2026,cloudrun2026}.

\begin{table}[H]
    \centering
    \caption{Công nghệ và vai trò cốt lõi}
    \label{tab:technology_core_roles}
    \renewcommand{\arraystretch}{1.15}
    \begin{tabularx}{\linewidth}{|p{0.22\linewidth}|p{0.30\linewidth}|X|}
        \hline
        \textbf{Nhóm công nghệ} & \textbf{Công nghệ} & \textbf{Vai trò cốt lõi trong hệ thống} \\ \hline
        Giao diện & React, Redux Toolkit, React Router, Axios & Xây dựng giao diện, quản lý trạng thái phía client, điều hướng và giao tiếp REST API. \\ \hline
        Backend nghiệp vụ & Java, Spring Boot, Spring Security, JPA, RestClient & Cung cấp REST API, xác thực, phân quyền, transaction và giao tiếp với AI Service. \\ \hline
        Dữ liệu chuẩn & MySQL & Lưu trữ tài khoản, sản phẩm, tồn kho, giỏ hàng, khuyến mãi và đơn hàng. \\ \hline
        AI Service & Python, FastAPI, Pydantic, Amazon Bedrock & Phân tích truy vấn, điều phối xử lý AI, kiểm tra dữ liệu vào/ra và sinh phản hồi. \\ \hline
        Hội thoại & LangGraph, Redis & Điều phối workflow nhiều lượt, lưu checkpoint và trạng thái phiên ngoài container. \\ \hline
        Đa phương thức & YOLOv26, Pillow, Cohere Embed v4 & Phát hiện/crop vùng sản phẩm, tiền xử lý ảnh và tạo embedding văn bản--hình ảnh. \\ \hline
        Truy hồi & Qdrant, hybrid fusion, reranker & Lưu vector, lọc metadata, truy hồi ứng viên và hợp nhất/xếp hạng kết quả. \\ \hline
        Triển khai & Docker, Docker Compose, GCE, Cloud Run, Secret Manager & Đóng gói, kết nối, triển khai các dịch vụ và quản lý cấu hình bí mật. \\ \hline
    \end{tabularx}
\end{table}

\subsection*{Kết chương}

Chương này đã rút gọn các lựa chọn công nghệ thành năm nhóm: tiêu chí và tổng quan;
ứng dụng và dữ liệu nghiệp vụ; AI và xử lý đa phương thức; truy hồi, xếp hạng và
trạng thái; triển khai và tổng hợp quyết định. Cấu trúc này giữ lại các quyết định
ảnh hưởng trực tiếp đến hệ thống, còn nguyên lý lý thuyết, kiến trúc chi tiết và bằng
chứng vận hành được trình bày ở các chương tương ứng. Lựa chọn cuối cùng ưu tiên khả
năng tích hợp, tái lập và kiểm soát dữ liệu, đồng thời duy trì các fallback cần thiết
cho các dịch vụ AI phụ thuộc bên ngoài.
